\documentclass[a4paper]{article}
\usepackage{amsfonts}
\usepackage{amsmath}
\usepackage{amssymb}
\usepackage{graphicx}
\newcommand{\dint}{\displaystyle\int}

\begin{document}

\noindent \large Auteur: Karanjot Singh \\
\noindent \large Studierichting: Informatica \\
\noindent \large Oefeningenreeks 6A nr. 9.1 \\

\normalsize

Gegeven: drie opeenvolgende termen van een rekenkundige rij waarvan de som en het product bekend zijn.\\

Laat de drie termen zijn: $(x - v)$, $x$, $(x + v)$ \\

Dan geldt het stelsel: \\

$\left\{
\begin{array}{l}
(x - v) + x + (x + v) = 51 \\
(x - v) \cdot x \cdot (x + v) = 4641 \\
\end{array}
\right.$ \\

Bereken de som: \\

$(x - v) + x + (x + v) = 3x = 51 \Rightarrow x = 17$ \\

Substitueer $x = 17$ in het product: \\

$(x - v) \cdot x \cdot (x + v) = x(x^2 - v^2)$ \\

$\Rightarrow x(x^2 - v^2) = 4641$ \\

$\Rightarrow 17(17^2 - v^2) = 4641$ \\

$\Rightarrow 17(289 - v^2) = 4641$ \\

$\Rightarrow 289 - v^2 = \dfrac{4641}{17} = 273$ \\

$\Rightarrow v^2 = 289 - 273 = 16$ \\

$\Rightarrow v = \pm4$ \\

Dus de drie termen zijn: \\

Voor $v = 4$: $(13,\ 17,\ 21)$ \\

Voor $v = -4$: $(21,\ 17,\ 13)$ \\

\begin{thebibliography}{999}
	%\bibitem{a} Referentie
\end{thebibliography}

\end{document}
